\chapter{Критические эксперименты (построение)}

Том \emph{construction} не считается закрытым без явного списка проверок:
что должно \emph{сломаться}, если неверны $g$, SI-замыкание или $\symAlpha$.

\begin{itemize}
  \item \textbf{Реестр verify} --- согласованность алгебры и SI-строки (\texttt{verify\_principles.py}); это не замена физического опыта, а gate на внутреннюю согласованность.
  \item \textbf{Macro T с микро-M} --- гладкость, дисперсия $\symOmegaLower(k)$, фронт импульса, $c=1$ на binomial-карте (MODEL~\S4.6--4.7; манифест~\S7).
  \item \textbf{Открыто / decisive} --- NLSE-коэффициент \texttt{T6}, полные live-sim на FCC/hex для заявленных T-предсказаний; формулировать как предсказание + критерий провала.
\end{itemize}

Дальнейшие тома серии заканчиваются аналогичным блоком под свой класс опытов (см.\ \texttt{book/SERIES.md}).
